% ====================================================================
% results.tex (dissertation example)
% ====================================================================

\section{Results}
\label{sec:results}

\subsection{Test 1: sine wave input}

The results of the sine wave test are given in \cref{tab:enobResults} and
plotted in \cref{fig:enobVsOsr}. The table is built from the file written by
the simulation, so the figures in it are exactly those the model produced.

\begin{table}[htbp]
\centering
\pgfplotstabletypeset[
  col sep=comma,
  columns={osr,theory,sim,error},
  columns/osr/.style={column name={Oversampling ratio}, column type=r},
  columns/theory/.style={column name={Theory (bits)}, fixed, fixed zerofill, precision=2, column type=r},
  columns/sim/.style={column name={Simulated (bits)}, fixed, fixed zerofill, precision=2, column type=r},
  columns/error/.style={column name={Error (LSB rms)}, fixed, fixed zerofill, precision=3, column type=r},
  every head row/.style={before row=\toprule, after row=\midrule},
  every last row/.style={after row=\bottomrule},
]{assets/data/enob_vs_osr.csv}
\caption{Effective resolution for a sine wave input with \qty{0.5}{LSB} rms of
noise, compared with the prediction of \cref{eq:enobTheory}.}
\label{tab:enobResults}
\end{table}

With no averaging the model gave 10.97 bits against a predicted 11.00, which
confirms that the noise and the quantisation error combine as
\cref{eq:averagedError} assumes. The resolution then rose by one bit for each
factor of four in oversampling ratio, reaching 15.00 bits at a ratio of 256.
The largest difference from the prediction at any ratio was 0.03 bits.

% ====================================================================
% enobVsOsr.tex
% --------------------------------------------------------------------
% Effective resolution against oversampling ratio, plotted with
% pgfplots straight from the simulation's CSV output.
% ====================================================================

\begin{figure}[htbp]
\centering
\begin{tikzpicture}
\begin{semilogxaxis}[
  width=0.88\linewidth,
  height=7cm,
  grid=major,
  log basis x=2,
  xmin=0.7, xmax=362,
  ymin=10.5, ymax=15.5,
  xtick={1,4,16,64,256},
  xticklabels={1,4,16,64,256},
  xlabel={Oversampling ratio, $M$},
  ylabel={Effective resolution (bits)},
  legend pos=north west,
  legend cell align=left,
]
  \addplot[thick, linkblue] table[col sep=comma, x=osr, y=theory]
    {assets/data/enob_vs_osr.csv};
  \addlegendentry{Theory}

  \addplot[only marks, mark=*, mark size=2.5pt, red]
    table[col sep=comma, x=osr, y=sim] {assets/data/enob_vs_osr.csv};
  \addlegendentry{Simulated}
\end{semilogxaxis}
\end{tikzpicture}
\caption{Effective resolution against oversampling ratio for a sine wave
input with \qty{0.5}{LSB} rms of noise.}
\label{fig:enobVsOsr}
\end{figure}


\subsection{Test 2: static inputs}

\Cref{fig:dcError} shows the rms error for static inputs at each of the four
noise levels, with the full figures in \cref{app:results}. The three
cases behave quite differently.

% ====================================================================
% dcError.tex
% --------------------------------------------------------------------
% Residual error for static inputs at four noise levels, plotted with
% pgfplots from the simulation's CSV output. The dashed line is the
% prediction for 0.5 LSB of noise, sqrt((1/12 + 0.25)/M).
% ====================================================================

\begin{figure}[htbp]
\centering
\begin{tikzpicture}
\begin{loglogaxis}[
  width=0.88\linewidth,
  height=8cm,
  grid=major,
  log basis x=2,
  xmin=0.7, xmax=362,
  ymin=0.02, ymax=2,
  xtick={1,4,16,64,256},
  xticklabels={1,4,16,64,256},
  ytick={0.02,0.05,0.1,0.2,0.5,1,2},
  yticklabels={0.02,0.05,0.1,0.2,0.5,1,2},
  xlabel={Oversampling ratio, $M$},
  ylabel={Error (LSB rms)},
  legend pos=south west,
  legend cell align=left,
  legend style={font=\footnotesize},
]
  \addplot[thick, mark=square*, black]
    table[col sep=comma, x=osr, y=n000] {assets/data/dc_error.csv};
  \addlegendentry{No noise}

  \addplot[thick, mark=triangle*, orange]
    table[col sep=comma, x=osr, y=n025] {assets/data/dc_error.csv};
  \addlegendentry{\qty{0.25}{LSB}}

  \addplot[thick, mark=*, linkblue]
    table[col sep=comma, x=osr, y=n050] {assets/data/dc_error.csv};
  \addlegendentry{\qty{0.5}{LSB}}

  \addplot[thick, mark=diamond*, red]
    table[col sep=comma, x=osr, y=n100] {assets/data/dc_error.csv};
  \addlegendentry{\qty{1.0}{LSB}}

  \addplot[dashed, gray, domain=1:256, samples=2] {sqrt((1/12 + 0.25)/x)};
  \addlegendentry{Theory for \qty{0.5}{LSB}}
\end{loglogaxis}
\end{tikzpicture}
\caption{Error in the averaged result for static inputs, at four levels of
input noise.}
\label{fig:dcError}
\end{figure}


With no noise the error stayed between 0.28 and \qty{0.29}{LSB} at every
ratio. This is the value of $q/\sqrt{12}$ from \cref{eq:quantisationNoise},
which means averaging 256 readings gave the same result as taking one.

With \qty{0.5}{LSB} and \qty{1.0}{LSB} of noise the error fell in proportion
to $1/\sqrt{M}$ across the whole range, following \cref{eq:averagedError}. At
a ratio of 256 the errors were \qty{0.036}{LSB} and \qty{0.067}{LSB}, which
correspond to 15.0 and 14.1 effective bits.

With \qty{0.25}{LSB} of noise the error fell at first, from
\qty{0.39}{LSB} to \qty{0.20}{LSB} at a ratio of 4, but then levelled out. At
a ratio of 256 it was \qty{0.070}{LSB}, about three times the
\qty{0.024}{LSB} that \cref{eq:averagedError} predicts, and little better
than at a ratio of 64.
